% =====================================================================
\chapter{Resultados}\label{cap:resultados}
% =====================================================================

\section{Consumo de líquidos y peso corporal}\label{sec:res-consumo}

\subsection{Consumo de líquidos durante el tratamiento}\label{sec:res-consumo-tratamiento}

El consumo de líquidos (agua y solución endulzada, normalizado al peso corporal) se analizó mediante ANOVA de medidas repetidas de cuatro vías, con la edad de exposición y el tratamiento (sacarosa o stevia) como factores entre sujetos y el tipo de fluido (agua o solución endulzada) y el tiempo como factores intrasujeto. Se detectó una interacción significativa edad $\times$ tratamiento $\times$ fluido $\times$ tiempo ($F(2{,}51, 60{,}21) = 7{,}97$; $p = 0{,}0003$), lo que indica que el patrón temporal de consumo de agua frente a solución endulzada difirió tanto según la edad de exposición como según el tipo de edulcorante; la interacción edad $\times$ tratamiento $\times$ tiempo también fue significativa ($F(2{,}89, 69{,}27) = 6{,}98$; $p = 0{,}0004$) (Figura~\ref{fig:consumo-liquidos}A,B).

Al promediar el consumo a lo largo de todo el período experimental, la interacción edad $\times$ tratamiento $\times$ fluido fue significativa ($F(1,24) = 11{,}99$; $p = 0{,}0020$): los animales juveniles consumieron, normalizado al peso corporal, considerablemente más solución endulzada que los adultos, tanto de stevia (2,8 veces más) como de sacarosa (2,4 veces más) (Tabla~\ref{tab:consumo-liquidos}; Figura~\ref{fig:consumo-liquidos-promedio}). Sin embargo, la relación entre el consumo de la solución endulzada y el de agua difirió según el edulcorante: el consumo de sacarosa se mantuvo consistentemente mayor que el de agua a lo largo de todo el período de exposición en ambas edades, mientras que el de stevia mostró una modulación temporal más marcada, particularmente en los animales juveniles, en quienes la stevia superó al agua entre PD27 y PD41 pero no a partir de PD43. En conjunto, estos resultados indican que la edad de exposición influye tanto en la magnitud del consumo normalizado de la bebida endulzada como en su patrón temporal, con la sacarosa manteniendo una predominancia más persistente sobre el agua que la stevia en ambas edades.

Para cuantificar directamente la preferencia por la solución endulzada, se calculó el índice de preferencia como el cociente entre el consumo de la solución endulzada y el consumo total de líquido (solución endulzada + agua), a partir de las medias de la Tabla~\ref{tab:consumo-liquidos} (no se dispone de los datos crudos por caja para este cálculo). El índice fue de 0,957 en el grupo sacarosa juvenil y de 0,950 en el adulto, y de 0,768 en el grupo stevia juvenil y de 0,782 en el adulto, lo que indica que la preferencia por cada solución endulzada fue similar entre ambas edades pese a la marcada diferencia de edad en el consumo normalizado al peso corporal descripta más arriba. Esta diferencia por edad, de hecho, no fue exclusiva de las soluciones endulzadas: también se observó en el consumo de agua de los animales control, mayor en el grupo juvenil que en el adulto (0,218 frente a 0,082\,mL/g de peso/día; Tabla~\ref{tab:consumo-liquidos}), lo que sugiere que refleja una diferencia general en el consumo de líquido normalizado al peso entre ventanas etarias, y no un efecto específico del tratamiento. A modo de referencia, y sin que esto implique una comparación formal, el consumo de sacarosa representó un aporte calórico estimado de 0,215\,kcal/g de peso/día en el grupo juvenil y de 0,092\,kcal/g de peso/día en el adulto ($0{,}1$\,g de sacarosa/mL $\times$ 4\,kcal/g), frente a un aporte del alimento balanceado \confirmar{kcal/g del alimento balanceado, para completar la comparación}.

\begin{table}[htbp]
\centering
\caption[Consumo de líquidos durante el tratamiento]{Consumo de líquidos durante el tratamiento, promediado a lo largo del período experimental. Valores expresados como media $\pm$ SEM (mL/g de peso corporal/día); entre paréntesis, número de cajas. $p$: comparación juvenil vs.\ adulto dentro de cada condición (prueba de Šídák).}
\label{tab:consumo-liquidos}
\footnotesize
\begin{tabularx}{\linewidth}{@{}l C C c@{}}
\toprule
\textbf{Condición} & \textbf{Juvenil} & \textbf{Adulto} & \textbf{$p$} \\
\midrule
Control -- Agua       & $0{,}218 \pm 0{,}013$ (7) & $0{,}082 \pm 0{,}001$ (7) & \textbf{$<$\,0,0001} \\
Sacarosa -- Agua      & $0{,}024 \pm 0{,}003$ (7) & $0{,}012 \pm 0{,}002$ (7) & \textbf{0,031} \\
Sacarosa -- Sacarosa  & $0{,}538 \pm 0{,}044$ (7) & $0{,}229 \pm 0{,}008$ (7) & \textbf{$<$\,0,0001} \\
Stevia -- Agua        & $0{,}073 \pm 0{,}007$ (7) & $0{,}024 \pm 0{,}005$ (7) & \textbf{0,0008} \\
Stevia -- Stevia      & $0{,}241 \pm 0{,}026$ (7) & $0{,}086 \pm 0{,}002$ (7) & \textbf{0,0003} \\
\bottomrule
\end{tabularx}
\end{table}

\begin{figure}[htbp]
  \centering
  \begin{subfigure}[b]{0.48\linewidth}
    \centering
    \incluirfigura[width=\linewidth]{Fig5_consumo_liquidos_juvenil.pdf}
    \caption{Juvenil}
  \end{subfigure}
  \hfill
  \begin{subfigure}[b]{0.48\linewidth}
    \centering
    \incluirfigura[width=\linewidth]{Fig5_consumo_liquidos_adulto.pdf}
    \caption{Adulto}
  \end{subfigure}
  \caption[Curso temporal del consumo de líquidos durante el tratamiento]{Curso temporal del consumo de líquidos durante el tratamiento. (A)~Grupo juvenil. (B)~Grupo adulto. Eje horizontal: día postnatal (PD); eje vertical: consumo de líquido normalizado al peso corporal (mL/g de peso corporal/día). Cada curva corresponde a una condición y a un fluido (agua o solución endulzada). Datos expresados como media $\pm$ SEM. $n = 7$ cajas por grupo.}
  \label{fig:consumo-liquidos}
\end{figure}

\begin{figure}[htbp]
  \centering
  \incluirfigura[width=0.9\linewidth]{Fig5_consumo_liquidos_promedio.pdf}
  \caption[Consumo de líquidos promediado durante el tratamiento]{Consumo de líquidos promediado a lo largo del período experimental, por condición y edad. Datos expresados como media $\pm$ SEM. Los asteriscos indican diferencias entre la exposición juvenil y la adulta dentro de cada condición (prueba de Šídák; $^{*}\,p < 0{,}05$; $^{***}\,p < 0{,}001$; $^{****}\,p < 0{,}0001$). $n = 7$ cajas por grupo.}
  \label{fig:consumo-liquidos-promedio}
\end{figure}

\subsection{Consumo de agua durante el lavado}\label{sec:res-consumo-lavado}

Durante el período de lavado posterior al tratamiento, en el que todos los animales tuvieron acceso exclusivo a agua, el consumo de agua (expresado como porcentaje del grupo control) se analizó mediante ANOVA de medidas repetidas de tres vías, con la edad y el tratamiento previo como factores entre sujetos y el tiempo como factor intrasujeto (corrección de Geisser-Greenhouse). Se observó un efecto principal del tiempo ($F(4{,}59, 110{,}09) = 17{,}08$; $p < 0{,}0001$), así como interacciones significativas edad $\times$ tiempo ($F(4{,}59, 110{,}09) = 9{,}03$; $p < 0{,}0001$), tratamiento $\times$ tiempo ($F(4{,}59, 110{,}09) = 3{,}29$; $p = 0{,}010$) y, de forma más relevante, edad $\times$ tratamiento $\times$ tiempo ($F(4{,}59, 110{,}09) = 7{,}58$; $p < 0{,}0001$), lo que indica que el patrón temporal del efecto del tratamiento previo sobre el consumo de agua durante el lavado difirió según la edad de exposición. También se observó una interacción edad $\times$ tratamiento sobre el valor promediado ($F(1,24) = 6{,}30$; $p = 0{,}019$), sin efectos principales aislados de la edad ($F(1,24) = 0{,}01$; $p = 0{,}919$) ni del tratamiento previo ($F(1,24) = 0{,}14$; $p = 0{,}714$).

Promediado a lo largo del lavado, el consumo de agua fue de $92{,}32 \pm 5{,}58$\,\% y $105{,}67 \pm 2{,}97$\,\% del control en los grupos juveniles previamente expuestos a sacarosa y a stevia, respectivamente, y de $103{,}48 \pm 5{,}44$\,\% y $93{,}57 \pm 4{,}04$\,\% en los correspondientes grupos adultos. Las comparaciones \textit{post hoc} de Šídák mostraron que, en los animales juveniles, los previamente expuestos a stevia consumieron más agua que los previamente expuestos a sacarosa en D11 ($p = 0{,}030$) y en D25 ($p = 0{,}005$), sin diferencias entre tratamientos en ningún día en los animales adultos. Al examinar el efecto de la edad dentro de cada tratamiento previo, los animales juveniles previamente expuestos a stevia mostraron valores mayores que los adultos en D11 ($p = 0{,}014$), D13 ($p = 0{,}0016$), D15 ($p = 0{,}0003$), D17 ($p = 0{,}0018$), D21 ($p = 0{,}0075$) y D23 ($p = 0{,}031$), mientras que no se detectaron diferencias entre edades en ningún día dentro del grupo previamente expuesto a sacarosa. Estos resultados indican una divergencia marcada, dependiente de la edad, en la respuesta temporal al consumo de agua tras la retirada de la stevia, que se hizo evidente principalmente durante la segunda mitad del período de lavado (Figura~\ref{fig:agua-lavado}).

\begin{figure}[htbp]
  \centering
  \incluirfigura[width=0.85\linewidth]{Fig5_agua_lavado.pdf}
  \caption[Consumo de agua durante el lavado]{Consumo de agua durante el período de lavado posterior al tratamiento, expresado como porcentaje del grupo control. Datos expresados como media $\pm$ SEM, según la edad y el tratamiento previo.}
  \label{fig:agua-lavado}
\end{figure}

\subsection{Consumo de alimento}\label{sec:res-consumo-chow}

El consumo de alimento balanceado, normalizado al peso corporal y promediado a lo largo del período experimental, se analizó mediante ANOVA de dos vías (edad $\times$ tratamiento), que reveló efectos significativos de la edad ($F(1,21) = 2363{,}85$; $p < 0{,}0001$) y del tratamiento ($F(2,21) = 57{,}37$; $p < 0{,}0001$), así como una interacción edad $\times$ tratamiento significativa ($F(2,21) = 3{,}86$; $p = 0{,}037$). Las comparaciones \textit{post hoc} de Šídák mostraron que el consumo de alimento fue mayor en los animales juveniles que en los adultos dentro de los tres tratamientos (control, sacarosa y stevia; $p < 0{,}0001$ en los tres casos). Dentro del grupo juvenil, el consumo fue menor en el grupo sacarosa que en el control ($p < 0{,}0001$) y que en el grupo stevia ($p < 0{,}0001$), sin diferencias entre control y stevia. Dentro del grupo adulto se observó una reducción similar en el grupo sacarosa respecto del control ($p = 0{,}0005$) y de la stevia ($p < 0{,}0001$), mientras que la stevia mostró un consumo mayor que el control ($p = 0{,}034$). Este patrón fue consistente con el curso temporal del consumo durante el tratamiento, en el que el efecto del tratamiento, el del tiempo y su interacción fueron significativos tanto en el grupo juvenil (tratamiento: $F(2,11) = 26{,}88$, $p = 0{,}0002$; tiempo: $F(2{,}34, 21{,}07) = 203{,}97$, $p < 0{,}0001$; interacción: $F(4{,}68, 21{,}07) = 11{,}54$, $p < 0{,}0001$) como en el adulto (tratamiento: $F(2,12) = 33{,}99$, $p < 0{,}0001$; tiempo: $F(1{,}51, 18{,}10) = 91{,}04$, $p < 0{,}0001$; interacción: $F(3{,}02, 18{,}10) = 14{,}99$, $p < 0{,}0001$) (Figura~\ref{fig:chow-juvenil}).

\begin{figure}[htbp]
  \centering
  \incluirfigura[width=0.7\linewidth]{Fig5_chow_juvenil.pdf}
  \caption[Curso temporal del consumo de alimento, grupo juvenil]{Curso temporal del consumo de alimento (normalizado al peso corporal) durante el tratamiento, grupo juvenil. \pendiente{Agregar gráfico equivalente del grupo adulto cuando esté disponible.} Datos expresados como media $\pm$ SEM.}
  \label{fig:chow-juvenil}
\end{figure}

\subsection{Peso corporal}\label{sec:res-consumo-peso}
\pendiente{redactar cuando estén los datos de peso corporal.}

\FloatBarrier
\section{Conducta tipo ansiosa (OF)}\label{sec:res-of}

El OF se analizó con las variables registradas por \anymaze{} (\S\ref{sec:of}): la distancia total recorrida; el número de entradas y la distancia recorrida en la zona central (zona 1) y en las zonas centrales combinadas (zonas 1+2); el número de zonas visitadas; y la conducta de \textit{rearing}, como número de eventos y tiempo. Cada variable se analizó mediante un ANOVA de dos vías (edad $\times$ tratamiento) seguido del post hoc de Tukey, que en este caso compara los tres tratamientos dentro de cada edad. La Tabla~\ref{tab:of-grupos} presenta las medias y el número de animales por grupo, la Tabla~\ref{tab:of-anova} resume los ANOVA y la Figura~\ref{fig:of} muestra los resultados. El número de animales difiere entre variables y entre grupos, y los grupos expuestos a stevia (6 o 7 animales) son los de menor tamaño frente a los grupos control y sacarosa (10 a 18 animales). \confirmar{por qué el número de animales difiere entre variables registradas en una misma sesión (p.\,ej., adulto-control: 11 animales en la distancia total, 17 en la distancia en la zona 1 y 10 en las entradas a las zonas 1+2)}.

\begin{table}[htbp]
\centering
\caption[Variables del OF por grupo experimental]{Variables del OF por grupo experimental. Valores expresados como media $\pm$ SEM; entre paréntesis, número de animales.}
\label{tab:of-grupos}
\tiny
\setlength{\tabcolsep}{2pt}
\begin{tabularx}{\linewidth}{@{}l C C C C C C@{}}
\toprule
& \multicolumn{3}{c}{\textbf{Juvenil}} & \multicolumn{3}{c}{\textbf{Adulto}} \\
\cmidrule(lr){2-4}\cmidrule(l){5-7}
\textbf{Variable} & \textbf{Control} & \textbf{Sacarosa} & \textbf{Stevia} & \textbf{Control} & \textbf{Sacarosa} & \textbf{Stevia} \\
\midrule
Distancia total (m) & $16{,}9 \pm 1{,}2$ (17) & $18{,}3 \pm 1{,}2$ (15) & $17{,}1 \pm 1{,}0$ (6) & $17{,}0 \pm 1{,}9$ (11) & $19{,}1 \pm 1{,}5$ (16) & $18{,}2 \pm 0{,}7$ (7) \\
Entradas a zona 1 & $3{,}8 \pm 0{,}9$ (16) & $2{,}9 \pm 0{,}5$ (15) & $4{,}2 \pm 1{,}8$ (6) & $0{,}7 \pm 0{,}3$ (17) & $1{,}2 \pm 0{,}3$ (17) & $4{,}4 \pm 1{,}2$ (7) \\
Distancia en zona 1 (m) & $0{,}31 \pm 0{,}07$ (18) & $0{,}28 \pm 0{,}05$ (15) & $0{,}39 \pm 0{,}16$ (6) & $0{,}09 \pm 0{,}03$ (17) & $0{,}17 \pm 0{,}05$ (17) & $0{,}37 \pm 0{,}09$ (7) \\
Entradas a zonas 1+2 & $13{,}9 \pm 1{,}0$ (15) & $7{,}1 \pm 0{,}8$ (14) & $33{,}5 \pm 3{,}2$ (6) & $13{,}3 \pm 2{,}6$ (10) & $12{,}6 \pm 1{,}6$ (11) & $37{,}0 \pm 2{,}6$ (7) \\
Distancia en zonas 1+2 (m) & $4{,}02 \pm 0{,}65$ (16) & $1{,}70 \pm 0{,}35$ (15) & $6{,}69 \pm 0{,}57$ (6) & $1{,}71 \pm 0{,}34$ (15) & $1{,}95 \pm 0{,}30$ (16) & $6{,}55 \pm 0{,}41$ (7) \\
Zonas visitadas (n.º) & $37{,}1 \pm 5{,}9$ (17) & $18{,}0 \pm 3{,}9$ (12) & $75{,}7 \pm 5{,}6$ (6) & $24{,}9 \pm 4{,}1$ (14) & $30{,}3 \pm 4{,}2$ (13) & $73{,}3 \pm 6{,}3$ (7) \\
\textit{Rearing}, n.º de eventos & $21{,}9 \pm 1{,}4$ (18) & $11{,}8 \pm 1{,}3$ (15) & $8{,}5 \pm 1{,}2$ (6) & $18{,}8 \pm 1{,}5$ (17) & $19{,}9 \pm 2{,}2$ (16) & $12{,}3 \pm 1{,}2$ (7) \\
\textit{Rearing}, tiempo (s) & $36{,}5 \pm 2{,}6$ (18) & $20{,}8 \pm 2{,}8$ (14) & $11{,}4 \pm 2{,}1$ (6) & $41{,}7 \pm 4{,}9$ (17) & $43{,}1 \pm 5{,}9$ (16) & $22{,}4 \pm 3{,}0$ (7) \\
\bottomrule
\end{tabularx}
\end{table}

\begin{table}[htbp]
\centering
\caption[Resumen del ANOVA de dos vías para las variables del OF]{Resumen del ANOVA de dos vías (edad $\times$ tratamiento) para las variables del OF. En negrita, los efectos significativos ($p < 0{,}05$).}
\label{tab:of-anova}
\footnotesize
\setlength{\tabcolsep}{3pt}
\begin{tabularx}{\linewidth}{@{}l c c c c c c@{}}
\toprule
& \multicolumn{2}{c}{\textbf{Tratamiento}} & \multicolumn{2}{c}{\textbf{Edad}} & \multicolumn{2}{c}{\textbf{Interacción}} \\
\cmidrule(lr){2-3}\cmidrule(lr){4-5}\cmidrule(l){6-7}
\textbf{Variable} & \textbf{$F$ (gl)} & \textbf{$p$} & \textbf{$F$ (gl)} & \textbf{$p$} & \textbf{$F$ (gl)} & \textbf{$p$} \\
\midrule
Distancia total & 0,90 (2,66) & 0,41 & 0,31 (1,66) & 0,58 & 0,06 (2,66) & 0,94 \\
Entradas a zona 1 & 4,11 (2,72) & \textbf{0,020} & 6,20 (1,72) & \textbf{0,015} & 2,29 (2,72) & 0,11 \\
Distancia en zona 1 & 2,80 (2,74) & 0,067 & 3,85 (1,74) & 0,053 & 0,94 (2,74) & 0,39 \\
Entradas a zonas 1+2 & 91,79 (2,57) & \textbf{$<$0,0001} & 3,51 (1,57) & 0,066 & 1,88 (2,57) & 0,16 \\
Distancia en zonas 1+2 & 38,88 (2,69) & \textbf{$<$0,0001} & 3,15 (1,69) & 0,080 & 5,05 (2,69) & \textbf{0,0090} \\
Zonas visitadas & 37,01 (2,63) & \textbf{$<$0,0001} & 0,03 (1,63) & 0,87 & 3,28 (2,63) & \textbf{0,044} \\
\textit{Rearing}, n.º de eventos & 12,69 (2,73) & \textbf{$<$0,0001} & 3,54 (1,73) & 0,064 & 6,67 (2,73) & \textbf{0,0022} \\
\textit{Rearing}, tiempo & 8,81 (2,72) & \textbf{0,0004} & 10,12 (1,72) & \textbf{0,0022} & 2,24 (2,72) & 0,11 \\
\bottomrule
\end{tabularx}
\end{table}

\begin{figure}[htbp]
  \centering
  \begin{subfigure}[b]{0.32\linewidth}
    \centering
    \incluirfigura[width=\linewidth]{Fig5_OF_distancia_total.pdf}
    \caption{}
  \end{subfigure}
  \hfill
  \begin{subfigure}[b]{0.32\linewidth}
    \centering
    \incluirfigura[width=\linewidth]{Fig5_OF_entradas_zona1.pdf}
    \caption{}
  \end{subfigure}
  \hfill
  \begin{subfigure}[b]{0.32\linewidth}
    \centering
    \incluirfigura[width=\linewidth]{Fig5_OF_distancia_zona1.pdf}
    \caption{}
  \end{subfigure}

  \vspace{0.8em}
  \begin{subfigure}[b]{0.32\linewidth}
    \centering
    \incluirfigura[width=\linewidth]{Fig5_OF_entradas_zona12.pdf}
    \caption{}
  \end{subfigure}
  \hfill
  \begin{subfigure}[b]{0.32\linewidth}
    \centering
    \incluirfigura[width=\linewidth]{Fig5_OF_distancia_zona12.pdf}
    \caption{}
  \end{subfigure}
  \hfill
  \begin{subfigure}[b]{0.32\linewidth}
    \centering
    \incluirfigura[width=\linewidth]{Fig5_OF_zonas_visitadas.pdf}
    \caption{}
  \end{subfigure}

  \vspace{0.8em}
  \begin{subfigure}[b]{0.32\linewidth}
    \centering
    \incluirfigura[width=\linewidth]{Fig5_OF_rearing_numero.pdf}
    \caption{}
  \end{subfigure}
  \hfill
  \begin{subfigure}[b]{0.32\linewidth}
    \centering
    \incluirfigura[width=\linewidth]{Fig5_OF_rearing_tiempo.pdf}
    \caption{}
  \end{subfigure}
  \hfill
  \begin{subfigure}[b]{0.32\linewidth}
    \centering
    \incluirfigura[width=0.8\linewidth]{Fig5_leyenda_grupos.png}
  \end{subfigure}
  \caption[Variables del \textit{open field} (OF)]{Variables del \textit{open field} (OF) en ratas expuestas a agua (control), sacarosa o stevia durante la ventana juvenil o adulta. (A)~Distancia total recorrida. (B)~Entradas a la zona central (zona 1). (C)~Distancia recorrida en la zona 1. (D)~Entradas a las zonas centrales (zonas 1+2). (E)~Distancia recorrida en las zonas 1+2. (F)~Número de zonas visitadas. (G)~Número de eventos de \textit{rearing}. (H)~Tiempo de \textit{rearing}. Las barras indican la media $\pm$ SEM. Azul: control; rojo: sacarosa; verde: stevia. Los corchetes indican las comparaciones de Tukey significativas entre tratamientos dentro de cada edad ($^{*}$~$p < 0{,}05$; $^{**}$~$p < 0{,}01$; $^{***}$~$p < 0{,}001$; $^{****}$~$p < 0{,}0001$); no se indican las no significativas.}
  \label{fig:of}
\end{figure}

\subsection{Actividad locomotora}\label{sec:res-of-locomocion}

La distancia total recorrida no difirió entre grupos: no hubo efecto del tratamiento ($F_{2,66} = 0{,}90$; $p = 0{,}41$) ni de la edad ($F_{1,66} = 0{,}31$; $p = 0{,}58$), ni interacción ($F_{2,66} = 0{,}06$; $p = 0{,}94$), y ninguna comparación de Tukey resultó significativa (Figura~\ref{fig:of}A). Las medias de los seis grupos se ubicaron entre $16{,}9$ y $19{,}1$\,m en los 5\,min de la sesión. No se detectaron, por lo tanto, diferencias en la actividad locomotora general, aunque el menor tamaño de los grupos expuestos a stevia limita la capacidad de detectar diferencias pequeñas en ellos.

\subsection{Exploración de las zonas centrales y de la arena}\label{sec:res-of-centro}

En la zona central (zona 1), la distancia recorrida no mostró interacción edad $\times$ tratamiento ($F_{2,74} = 0{,}94$; $p = 0{,}39$), y los efectos principales del tratamiento ($F_{2,74} = 2{,}80$; $p = 0{,}067$) y de la edad ($F_{1,74} = 3{,}85$; $p = 0{,}053$) no alcanzaron la significancia estadística. En los adultos, los animales expuestos a stevia recorrieron más distancia en la zona 1 que los controles (diferencia control $-$ stevia $= -0{,}28$\,m; IC~95\,\%: $-0{,}54$ a $-0{,}02$; $p < 0{,}05$; Figura~\ref{fig:of}C); en los juveniles no hubo diferencias entre tratamientos. El número de entradas a la zona 1 mostró efectos principales del tratamiento ($F_{2,72} = 4{,}11$; $p = 0{,}020$) y de la edad ($F_{1,72} = 6{,}20$; $p = 0{,}015$), sin interacción ($F_{2,72} = 2{,}29$; $p = 0{,}11$): en los adultos, los animales expuestos a stevia ingresaron más veces a la zona 1 que los controles (diferencia control $-$ stevia $= -3{,}7$; IC~95\,\%: $-6{,}4$ a $-1{,}1$; $p < 0{,}01$) y que los expuestos a sacarosa (diferencia sacarosa $-$ stevia $= -3{,}2$; IC~95\,\%: $-5{,}8$ a $-0{,}6$; $p < 0{,}05$), y en los juveniles no hubo diferencias (Figura~\ref{fig:of}B).

Al considerar las zonas centrales en conjunto (zonas 1+2), el efecto del tratamiento fue marcado. El número de entradas dependió claramente del tratamiento ($F_{2,57} = 91{,}79$; $p < 0{,}0001$), sin interacción ($F_{2,57} = 1{,}88$; $p = 0{,}16$) ni efecto de la edad ($F_{1,57} = 3{,}51$; $p = 0{,}066$). En ambas edades, los animales expuestos a stevia ingresaron más veces a las zonas centrales que los controles y que los expuestos a sacarosa (Figura~\ref{fig:of}D). En los juveniles, la diferencia control $-$ stevia fue de $-19{,}6$ entradas (IC~95\,\%: $-26{,}1$ a $-13{,}0$; $p < 0{,}0001$) y la diferencia sacarosa $-$ stevia, de $-26{,}4$ (IC~95\,\%: $-33{,}1$ a $-19{,}8$; $p < 0{,}0001$); en los adultos, de $-23{,}7$ (IC~95\,\%: $-30{,}4$ a $-17{,}0$; $p < 0{,}0001$) y de $-24{,}4$ (IC~95\,\%: $-30{,}9$ a $-17{,}8$; $p < 0{,}0001$), respectivamente. Los animales expuestos a sacarosa presentaron menos entradas que los controles solo en los juveniles (diferencia control $-$ sacarosa $= 6{,}9$; IC~95\,\%: $1{,}8$ a $11{,}9$; $p < 0{,}01$), sin diferencia respecto del control en los adultos ($0{,}7$; IC~95\,\%: $-5{,}3$ a $6{,}6$).

La distancia recorrida en las zonas 1+2 mostró una interacción edad $\times$ tratamiento significativa ($F_{2,69} = 5{,}05$; $p = 0{,}009$) y un efecto principal del tratamiento ($F_{2,69} = 38{,}88$; $p < 0{,}0001$), sin efecto de la edad ($F_{1,69} = 3{,}15$; $p = 0{,}080$). En los juveniles, los animales expuestos a stevia recorrieron más distancia que los controles (diferencia control $-$ stevia $= -2{,}7$\,m; IC~95\,\%: $-4{,}6$ a $-0{,}8$; $p < 0{,}01$) y que los expuestos a sacarosa (diferencia sacarosa $-$ stevia $= -5{,}0$; IC~95\,\%: $-6{,}9$ a $-3{,}1$; $p < 0{,}0001$), y estos últimos recorrieron menos que los controles (diferencia control $-$ sacarosa $= 2{,}3$; IC~95\,\%: $0{,}9$ a $3{,}7$; $p < 0{,}001$). En los adultos, la stevia superó al control (diferencia control $-$ stevia $= -4{,}8$; IC~95\,\%: $-6{,}7$ a $-3{,}0$; $p < 0{,}0001$) y a la sacarosa (diferencia sacarosa $-$ stevia $= -4{,}6$; IC~95\,\%: $-6{,}4$ a $-2{,}8$; $p < 0{,}0001$), sin diferencia entre control y sacarosa ($-0{,}2$; IC~95\,\%: $-1{,}7$ a $1{,}2$) (Figura~\ref{fig:of}E).

El número de zonas visitadas mostró una interacción edad $\times$ tratamiento significativa ($F_{2,63} = 3{,}28$; $p = 0{,}044$) y un efecto principal del tratamiento ($F_{2,63} = 37{,}01$; $p < 0{,}0001$), sin efecto de la edad ($F_{1,63} = 0{,}03$; $p = 0{,}87$). En ambas edades, los animales expuestos a stevia visitaron más zonas que los controles y que los expuestos a sacarosa (Figura~\ref{fig:of}F): en los juveniles, la diferencia control $-$ stevia fue de $-38{,}6$ (IC~95\,\%: $-58{,}8$ a $-18{,}3$; $p < 0{,}0001$) y la diferencia sacarosa $-$ stevia, de $-57{,}7$ (IC~95\,\%: $-79{,}0$ a $-36{,}3$; $p < 0{,}0001$); en los adultos, de $-48{,}4$ (IC~95\,\%: $-68{,}1$ a $-28{,}6$; $p < 0{,}0001$) y de $-43{,}0$ (IC~95\,\%: $-63{,}0$ a $-23{,}0$; $p < 0{,}0001$). Los juveniles expuestos a sacarosa visitaron menos zonas que los controles (diferencia control $-$ sacarosa $= 19{,}1$; IC~95\,\%: $3{,}0$ a $35{,}2$; $p < 0{,}05$), sin diferencia en los adultos ($-5{,}4$; IC~95\,\%: $-21{,}8$ a $11{,}1$). Estos aumentos con stevia ocurrieron sin diferencias en la distancia total recorrida (\S\ref{sec:res-of-locomocion}).

\subsection{Conducta de \textit{rearing}}\label{sec:res-of-rearing}

El número de eventos de \textit{rearing} mostró una interacción edad $\times$ tratamiento significativa ($F_{2,73} = 6{,}67$; $p = 0{,}0022$) y un efecto principal del tratamiento ($F_{2,73} = 12{,}69$; $p < 0{,}0001$), sin efecto principal de la edad ($F_{1,73} = 3{,}54$; $p = 0{,}064$). En los juveniles, tanto los animales expuestos a sacarosa como los expuestos a stevia presentaron menos eventos de \textit{rearing} que los controles (diferencia control $-$ sacarosa $= 10{,}1$; IC~95\,\%: $4{,}9$ a $15{,}4$; $p < 0{,}0001$; diferencia control $-$ stevia $= 13{,}4$; IC~95\,\%: $6{,}4$ a $20{,}5$; $p < 0{,}0001$), sin diferencias entre ambos tratamientos (Figura~\ref{fig:of}G). En los adultos, el \textit{rearing} con sacarosa no difirió del control (diferencia control $-$ sacarosa $= -1{,}1$; IC~95\,\%: $-6{,}4$ a $4{,}1$) y fue menor con stevia que con sacarosa (diferencia sacarosa $-$ stevia $= 7{,}7$; IC~95\,\%: $0{,}8$ a $14{,}5$; $p < 0{,}05$), sin alcanzar la significancia respecto del control (diferencia control $-$ stevia $= 6{,}5$; IC~95\,\%: $-0{,}2$ a $13{,}3$).

El tiempo de \textit{rearing} mostró efectos principales del tratamiento ($F_{2,72} = 8{,}81$; $p = 0{,}0004$) y de la edad ($F_{1,72} = 10{,}12$; $p = 0{,}0022$), sin interacción significativa ($F_{2,72} = 2{,}24$; $p = 0{,}11$). En los juveniles, fue menor con sacarosa (diferencia control $-$ sacarosa $= 15{,}7$\,s; IC~95\,\%: $1{,}8$ a $29{,}5$; $p < 0{,}05$) y con stevia (diferencia control $-$ stevia $= 25{,}1$\,s; IC~95\,\%: $6{,}7$ a $43{,}4$; $p < 0{,}01$) que con el control, sin diferencia entre ambos tratamientos ($9{,}4$\,s; IC~95\,\%: $-9{,}6$ a $28{,}4$). En los adultos, fue menor con stevia que con el control (diferencia control $-$ stevia $= 19{,}3$\,s; IC~95\,\%: $1{,}8$ a $36{,}8$; $p < 0{,}05$) y que con sacarosa (diferencia sacarosa $-$ stevia $= 20{,}7$\,s; IC~95\,\%: $3{,}1$ a $38{,}3$; $p < 0{,}05$), sin diferencia entre control y sacarosa ($-1{,}4$\,s; IC~95\,\%: $-15{,}0$ a $12{,}1$) (Figura~\ref{fig:of}H).

\FloatBarrier
\subsection{Síntesis}\label{sec:res-of-sintesis}

En conjunto, la locomoción general no difirió entre grupos. Los animales expuestos a stevia presentaron más entradas, mayor distancia recorrida y más zonas visitadas en las zonas centrales 1+2 y en la arena que los controles y que los expuestos a sacarosa, en ambas edades, mientras que en la zona 1 la diferencia respecto del control apareció solo en los adultos. Los animales expuestos a sacarosa difirieron del control solo en la ventana juvenil, con menos entradas y menor distancia en las zonas 1+2, menos zonas visitadas y menos \textit{rearing}. El \textit{rearing} fue menor con stevia que con el control en los juveniles (en número y en tiempo) y, en los adultos, menor que con sacarosa (en número y en tiempo) y que con el control (solo en tiempo).

\FloatBarrier
\section{Memoria de reconocimiento de objetos (NOR)}\label{sec:res-nor}

La memoria de reconocimiento de objetos se evaluó mediante la tasa de exploración del objeto novedoso (nivel de azar 0,5) en tres demoras crecientes respecto de la adquisición (T1): T2 (2\,h, memoria a corto plazo), T3 (24\,h, memoria a largo plazo) y T4 (7 días, memoria remota), analizadas por separado mediante ANOVA de dos vías (edad de exposición $\times$ tratamiento), con el mismo esquema de comparaciones \textit{post hoc} empleado en la prueba SLR (\S\ref{sec:res-slr}). La Tabla~\ref{tab:nor-tasa} resume la tasa de exploración de cada grupo y su comparación con el nivel de azar, la Tabla~\ref{tab:nor-anova} resume los análisis de varianza y la Figura~\ref{fig:nor-tasa} muestra la distribución individual de los datos. En las tres demoras la interacción edad $\times$ tratamiento resultó significativa, por lo que los efectos principales se interpretan con cautela y el análisis se centra en las comparaciones \textit{post hoc} dentro de cada edad y de cada tratamiento.

\subsection{Validación de la prueba}\label{sec:res-nor-validacion}

En las tres demoras, los dos grupos control exploraron preferentemente el objeto novedoso por encima del nivel de azar, lo que confirma que la prueba detecta memoria de reconocimiento en ambas ventanas etarias. En el grupo juvenil, la tasa de exploración del control fue de 0,632 a las 2\,h ($t(8) = 10{,}31$; $p < 0{,}0001$), 0,716 a las 24\,h ($t(8) = 8{,}62$; $p < 0{,}0001$) y 0,677 a los 7 días ($t(8) = 15{,}66$; $p < 0{,}0001$); en el grupo adulto fue de 0,703 a las 2\,h ($t(9) = 10{,}56$; $p < 0{,}0001$), 0,653 a las 24\,h ($t(9) = 4{,}72$; $p = 0{,}0011$) y 0,662 a los 7 días ($t(9) = 4{,}61$; $p = 0{,}0013$) (Tabla~\ref{tab:nor-tasa}). Sobre esa línea de base, el efecto de cada edulcorante dependió de la ventana de exposición y de la demora, como lo indica la interacción edad $\times$ tratamiento, significativa en las tres demoras (T2: $F(2,49) = 10{,}13$, $p = 0{,}0002$; T3: $F(2,49) = 11{,}67$, $p < 0{,}0001$; T4: $F(2,49) = 7{,}82$, $p = 0{,}0011$; Tabla~\ref{tab:nor-anova}).

\subsection{Efecto de la sacarosa}\label{sec:res-nor-sacarosa}

El hallazgo más consistente de la prueba es que la exposición juvenil a sacarosa redujo la tasa de exploración del objeto novedoso respecto del control juvenil en las tres demoras (T2: diferencia control $-$ sacarosa $= 0{,}138$; IC~95\,\%: 0,047--0,228; $p = 0{,}0017$; T3: diferencia control $-$ sacarosa $= 0{,}150$; IC~95\,\%: 0,058--0,242; $p = 0{,}0008$; T4: diferencia control $-$ sacarosa $= 0{,}172$; IC~95\,\%: 0,065--0,280; $p = 0{,}0009$; Figura~\ref{fig:nor-tasa}). La sacarosa adulta, en cambio, no difirió de su control en ninguna de las tres demoras (T2: $p = 0{,}962$; T3: sin diferencias entre tratamientos, $p \geq 0{,}077$; T4: sin diferencia frente al control). La diferencia entre la exposición juvenil y la adulta dentro del grupo sacarosa fue significativa y de magnitud comparable en las tres demoras (T2: diferencia juvenil $-$ adulto $= -0{,}199$; $p < 0{,}0001$; T3: diferencia juvenil $-$ adulto $= -0{,}170$; $p < 0{,}0001$; T4: diferencia juvenil $-$ adulto $= -0{,}217$; $p < 0{,}0001$), mientras que el control no difirió entre edades en ninguna demora ($p = 0{,}059$, $p = 0{,}097$ y $p = 0{,}718$ en T2, T3 y T4, respectivamente).

\subsection{Efecto de la stevia}\label{sec:res-nor-stevia}

El efecto de la stevia, en contraste, no fue consistente entre demoras ni compartido entre edades. En el grupo juvenil, la stevia solo difirió del control a las 24\,h (diferencia control $-$ stevia $= 0{,}155$; IC~95\,\%: 0,063--0,248; $p = 0{,}0005$), sin diferencias significativas a las 2\,h ($p = 0{,}461$) ni a los 7 días (diferencia control $-$ stevia $= 0{,}097$; IC~95\,\%: $-0{,}011$--0,204; $p = 0{,}087$). En el grupo adulto el patrón se invirtió: la stevia solo difirió del control a las 2\,h (diferencia control $-$ stevia $= 0{,}155$; IC~95\,\%: 0,067--0,243; $p = 0{,}0003$), sin diferencias a las 24\,h (sin diferencias entre tratamientos, $p \geq 0{,}077$) ni, frente al control, a los 7 días (aunque sí respecto de la sacarosa: diferencia sacarosa $-$ stevia $= 0{,}113$; IC~95\,\%: 0,005--0,220; $p = 0{,}038$). Ninguno de los dos grupos stevia difirió entre edades en ninguna demora ($p \geq 0{,}296$ en todas las comparaciones).

\begin{table}[htbp]
\centering
\caption[Tasa de exploración en la prueba NOR y comparación con el azar]{Tasa de exploración en la prueba NOR a cada demora y comparación con el nivel de azar (0,5) mediante prueba $t$ de una muestra. Valores expresados como media $\pm$ SEM; entre paréntesis, número de animales.}
\label{tab:nor-tasa}
\tiny
\setlength{\tabcolsep}{3pt}
\begin{tabularx}{\linewidth}{@{}l C c c C c c C c c@{}}
\toprule
& \multicolumn{3}{c}{\textbf{T2 (2\,h)}} & \multicolumn{3}{c}{\textbf{T3 (24\,h)}} & \multicolumn{3}{c}{\textbf{T4 (7 días)}} \\
\cmidrule(lr){2-4}\cmidrule(lr){5-7}\cmidrule(l){8-10}
\textbf{Grupo} & \textbf{Tasa} & \textbf{$t$ (gl)} & \textbf{$p$} & \textbf{Tasa} & \textbf{$t$ (gl)} & \textbf{$p$} & \textbf{Tasa} & \textbf{$t$ (gl)} & \textbf{$p$} \\
\midrule
Juvenil -- Control  & $0{,}632 \pm 0{,}013$ (9)  & 10,31 (8) & \textbf{$<$0,0001} & $0{,}716 \pm 0{,}025$ (9) & 8,62 (8)  & \textbf{$<$0,0001} & $0{,}677 \pm 0{,}011$ (9) & 15,66 (8) & \textbf{$<$0,0001} \\
Juvenil -- Sacarosa & $0{,}494 \pm 0{,}040$ (9)  & $-$0,15 (8) & 0,885 & $0{,}566 \pm 0{,}032$ (9) & 2,05 (8)  & 0,074 & $0{,}505 \pm 0{,}024$ (9) & 0,22 (8) & 0,830 \\
Juvenil -- Stevia   & $0{,}587 \pm 0{,}023$ (9)  & 3,80 (8) & \textbf{0,0052} & $0{,}561 \pm 0{,}021$ (9) & 2,88 (8)  & \textbf{0,021} & $0{,}581 \pm 0{,}021$ (9) & 3,83 (8) & \textbf{0,0050} \\
Adulto -- Control   & $0{,}703 \pm 0{,}019$ (10) & 10,56 (9) & \textbf{$<$0,0001} & $0{,}653 \pm 0{,}032$ (10) & 4,72 (9) & \textbf{0,0011} & $0{,}662 \pm 0{,}035$ (10) & 4,61 (9) & \textbf{0,0013} \\
Adulto -- Sacarosa  & $0{,}693 \pm 0{,}024$ (9)  & 7,98 (8) & \textbf{$<$0,0001} & $0{,}736 \pm 0{,}021$ (9) & 11,40 (8) & \textbf{$<$0,0001} & $0{,}722 \pm 0{,}036$ (9) & 6,23 (8) & \textbf{0,0003} \\
Adulto -- Stevia    & $0{,}548 \pm 0{,}032$ (9)  & 1,50 (8) & 0,172 & $0{,}707 \pm 0{,}024$ (9) & 8,55 (8)  & \textbf{$<$0,0001} & $0{,}610 \pm 0{,}046$ (9) & 2,37 (8) & \textbf{0,045} \\
\bottomrule
\end{tabularx}
\end{table}

\begin{table}[htbp]
\centering
\caption[Resumen del ANOVA de dos vías para la tasa de exploración en la prueba NOR]{Resumen del ANOVA de dos vías (edad de exposición $\times$ tratamiento) para la tasa de exploración en cada demora de la prueba NOR. $\eta^2_p$: eta cuadrado parcial.}
\label{tab:nor-anova}
\footnotesize
\setlength{\tabcolsep}{3pt}
\begin{tabularx}{\linewidth}{@{}l c c C c c C c c C@{}}
\toprule
& \multicolumn{3}{c}{\textbf{T2}} & \multicolumn{3}{c}{\textbf{T3}} & \multicolumn{3}{c}{\textbf{T4}} \\
\cmidrule(lr){2-4}\cmidrule(lr){5-7}\cmidrule(l){8-10}
\textbf{Fuente} & \textbf{$F$ (gl)} & \textbf{$p$} & \textbf{$\eta^2_p$} & \textbf{$F$ (gl)} & \textbf{$p$} & \textbf{$\eta^2_p$} & \textbf{$F$ (gl)} & \textbf{$p$} & \textbf{$\eta^2_p$} \\
\midrule
Tratamiento & 7,89 (2,49) & \textbf{0,0011} & 0,244 & 1,87 (2,49) & 0,165 & 0,071 & 3,11 (2,49) & 0,054 & 0,113 \\
Edad        & 12,74 (1,49) & \textbf{0,0008} & 0,206 & 14,92 (1,49) & \textbf{0,0003} & 0,233 & 9,07 (1,49) & \textbf{0,0041} & 0,156 \\
Interacción & 10,13 (2,49) & \textbf{0,0002} & 0,293 & 11,67 (2,49) & \textbf{$<$0,0001} & 0,323 & 7,82 (2,49) & \textbf{0,0011} & 0,242 \\
\bottomrule
\end{tabularx}
\end{table}

\begin{figure}[htbp]
  \centering
  \begin{subfigure}[b]{0.32\linewidth}
    \centering
    \incluirfigura[width=\linewidth]{Fig5_NOR_T2.pdf}
    \caption{T2 (2\,h)}
  \end{subfigure}
  \hfill
  \begin{subfigure}[b]{0.32\linewidth}
    \centering
    \incluirfigura[width=\linewidth]{Fig5_NOR_T3.pdf}
    \caption{T3 (24\,h)}
  \end{subfigure}
  \hfill
  \begin{subfigure}[b]{0.32\linewidth}
    \centering
    \incluirfigura[width=\linewidth]{Fig5_NOR_T4.pdf}
    \caption{T4 (7 días)}
  \end{subfigure}
  \vspace{0.6em}
  \incluirfigura[width=0.2\linewidth]{Fig5_leyenda_grupos.png}
  \caption[Tasa de exploración en la prueba NOR]{Tasa de exploración en la prueba NOR. (A)~T2 (2\,h). (B)~T3 (24\,h). (C)~T4 (7 días). Las barras indican media $\pm$ SEM; cada punto representa un animal. Azul: control; rojo: sacarosa; verde: stevia (paleta común a todas las figuras de esta tesis). Se indican los resultados de las comparaciones de Tukey entre grupos ($^{*}\,p < 0{,}05$; $^{**}\,p < 0{,}01$; $^{***}\,p < 0{,}001$; $^{****}\,p < 0{,}0001$). $n = 9$--$10$ animales por grupo.}
  \label{fig:nor-tasa}
\end{figure}

\begin{figure}[htbp]
  \centering
  \incluirfigura[width=\linewidth]{Fig5_estimacion_Tukey_NOR.png}
  \caption[Diferencias de medias de la tasa de exploración entre grupos en la prueba NOR (prueba de Tukey)]{Diferencias de medias de la tasa de exploración entre grupos (prueba de Tukey) con su intervalo de confianza del 95\,\%, para las demoras T2 (A), T3 (B) y T4 (C). Cada fila corresponde a una comparación; una diferencia positiva indica una tasa mayor en el primer grupo de la comparación. En rojo se indican las comparaciones significativas ($p < 0{,}05$). La línea punteada indica ausencia de diferencia.}
  \label{fig:nor-tukey}
\end{figure}

\FloatBarrier
\subsection{Síntesis}\label{sec:res-nor-sintesis}

En conjunto, estos resultados señalan que el efecto de la sacarosa sobre la memoria de reconocimiento es robusto frente a la demora pero circunscripto a la ventana juvenil, mientras que el efecto de la stevia es más acotado y depende tanto de la edad de exposición como del intervalo de retención evaluado, sin un patrón único que permita ubicarla de forma estable entre el control y la sacarosa. La Figura~\ref{fig:nor-tukey} resume las diferencias entre grupos junto con sus intervalos de confianza.

% ---------------------------------------------------------------------
\FloatBarrier
\section{Separación de patrones espaciales (SLR)}\label{sec:res-slr}

La memoria espacial y la capacidad de separación de patrones se evaluaron mediante la prueba SLR en sus dos configuraciones, d-SLR y s-SLR (Figura~\ref{fig:esquema-slr}), que se analizaron por separado. El desempeño se cuantificó con los dos índices definidos en \S\ref{sec:nor} y \S\ref{sec:slr} a partir de los mismos tiempos de exploración activa: el índice de discriminación (DI, nivel de azar 0) y la tasa de exploración (nivel de azar 0,5), relacionados por la transformación lineal $\mathrm{DI} = 2\cdot\text{tasa} - 1$. Por tratarse de una transformación lineal, ambos índices producen resultados estadísticos idénticos (mismos $F$, $t$ y $p$); se presentan los dos a continuación porque difieren en su escala e interpretación. La Tabla~\ref{tab:slr-di} y la Figura~\ref{fig:slr-di} resumen el DI de cada grupo y su comparación con el nivel de azar, la Tabla~\ref{tab:slr-tasa} y la Figura~\ref{fig:slr-tasa} presentan la tasa de exploración equivalente, y la Tabla~\ref{tab:slr-anova} resume los análisis de varianza.

\subsection{Exploración durante las fases de muestreo y de prueba}\label{sec:res-slr-exploracion}

La exploración total (tiempo de exploración activa de los objetos, puntuado manualmente) del día 1 y del día 2 se muestra en la Figura~\ref{fig:slr-exploracion}; funciona como control de la codificación y no como desenlace (\S\ref{sec:slr}). La exploración disminuyó del día 1 al día 2 en todos los grupos (entre 44\,\% y 59\,\%). En el día 1 no fue equivalente entre grupos: fue mayor en los juveniles expuestos a sacarosa ($54{,}5 \pm 3{,}6$\,s) y en los adultos expuestos a stevia ($55{,}8 \pm 5{,}5$\,s) que en el resto ($30{,}6$ a $41{,}7$\,s), con una interacción edad $\times$ tratamiento significativa ($F_{2,111} = 6{,}72$; $p = 0{,}002$; ANOVA de dos vías sobre el logaritmo del tiempo). Cada observación corresponde a un animal (los grupos de d-SLR y de s-SLR están formados por animales distintos; Anexo~\ref{anx:muestra}), y las dos configuraciones se combinaron en este análisis. La diferencia se tiene presente al interpretar los índices de discriminación de las secciones siguientes.

\begin{figure}[htbp]
  \centering
  \incluirfigura[width=0.85\linewidth]{Fig5_exploracion_SLR.pdf}
  \caption[Exploración total en el día 1 y en el día 2 por grupo (SLR)]{Exploración total (tiempo de exploración activa de los objetos, en segundos) en el día 1 (muestreo, barras llenas) y en el día 2 (prueba, barras rayadas) para cada grupo, con exposición juvenil (izquierda) y adulta (derecha). Las barras indican la media $\pm$ SEM y los puntos, las sesiones individuales, con las sesiones de d-SLR y s-SLR combinadas. Azul: control; rojo: sacarosa; verde: stevia. Número de sesiones por grupo, igual en ambos días: juvenil, control 19, sacarosa 18 y stevia 20; adulto, 20 en cada grupo.}
  \label{fig:slr-exploracion}
\end{figure}

\subsection{Validación de la prueba}\label{sec:res-slr-validacion}

En ambas configuraciones, solo los grupos control discriminaron la posición novedosa por encima del nivel de azar, tanto en la exposición juvenil (d-SLR: $p = 0{,}001$; s-SLR: $p < 0{,}001$) como en la adulta (d-SLR: $p = 0{,}003$; s-SLR: $p = 0{,}007$), lo que confirma que la prueba detecta separación de patrones espaciales en ambas ventanas etarias y bajo ambas demandas de carga cognitiva (Tabla~\ref{tab:slr-di}). Ninguno de los grupos expuestos a sacarosa o a stevia difirió significativamente del azar en ninguna configuración (Tabla~\ref{tab:slr-di}; Figura~\ref{fig:slr-di}). El ANOVA de dos vías reveló, en ambas configuraciones, un efecto significativo del tratamiento (d-SLR: $F(2,52) = 4{,}58$; $p = 0{,}015$; $\eta^2_p = 0{,}15$; s-SLR: $F(2,53) = 6{,}67$; $p = 0{,}003$; $\eta^2_p = 0{,}20$), sin efecto de la edad de exposición (d-SLR: $F(1,52) = 0{,}22$; $p = 0{,}642$; s-SLR: $F(1,53) = 0{,}08$; $p = 0{,}779$) ni de la interacción entre ambos factores (d-SLR: $F(2,52) = 2{,}04$; $p = 0{,}140$; s-SLR: $F(2,53) = 0{,}53$; $p = 0{,}591$) (Tabla~\ref{tab:slr-anova}).

\subsection{Efecto de la sacarosa}\label{sec:res-slr-sacarosa}

En la configuración d-SLR, el DI de los animales expuestos a sacarosa fue significativamente menor que el de los controles en el grupo juvenil (diferencia control $-$ sacarosa $= 0{,}324$; IC~95\,\%: 0,071--0,577; $p = 0{,}009$), sin diferencias en el grupo adulto ($p > 0{,}50$ en todas las comparaciones); la comparación entre edades dentro del grupo sacarosa mostró solo una tendencia no significativa ($p = 0{,}067$). En la configuración s-SLR, en cambio, la sacarosa redujo el DI respecto del control en ambas edades (juvenil: diferencia control $-$ sacarosa $= 0{,}322$; IC~95\,\%: 0,011--0,633; $p = 0{,}041$; adulto: diferencia control $-$ sacarosa $= 0{,}326$; IC~95\,\%: 0,023--0,629; $p = 0{,}032$), sin diferencias entre edades dentro del tratamiento. Los residuos del modelo de s-SLR se apartaron de la normalidad (Shapiro-Wilk, $p = 0{,}006$), por lo que estas diferencias se contrastaron también mediante la prueba de Mann-Whitney, que las confirmó en ambas edades (juvenil: $p = 0{,}002$; adulto: $p = 0{,}031$); la homogeneidad de varianzas se mantuvo (Brown-Forsythe, $p = 0{,}16$), al igual que en d-SLR (Shapiro-Wilk, $p = 0{,}48$; Brown-Forsythe, $p = 0{,}30$).

\subsection{Efecto de la stevia}\label{sec:res-slr-stevia}

La stevia no difirió del control en ninguna configuración ni edad (d-SLR, juvenil: $p = 0{,}870$; d-SLR, adulto: $p > 0{,}50$; s-SLR, juvenil: $p = 0{,}971$; s-SLR, adulto: $p = 0{,}293$). En la configuración d-SLR, el DI del grupo stevia juvenil fue mayor que el del grupo sacarosa juvenil (diferencia stevia $-$ sacarosa $= 0{,}272$; IC~95\,\%: 0,026--0,519; $p = 0{,}027$), mientras que esta misma comparación en s-SLR no alcanzó significancia estadística ($p = 0{,}069$). Ningún grupo stevia difirió entre edades, y ninguno superó significativamente el nivel de azar en ninguna configuración (Tabla~\ref{tab:slr-di}).

\begin{table}[htbp]
\centering
\caption[Índice de discriminación (DI) en la prueba SLR]{Índice de discriminación (DI) en la prueba SLR y comparación con el nivel de azar ($\mathrm{DI} = 0$) mediante prueba $t$ de una muestra. Valores expresados como media $\pm$ SEM; entre paréntesis, número de animales.}
\label{tab:slr-di}
\footnotesize
\setlength{\tabcolsep}{4pt}
\begin{tabularx}{\linewidth}{@{}l C c c C c c@{}}
\toprule
& \multicolumn{3}{c}{\textbf{d-SLR}} & \multicolumn{3}{c}{\textbf{s-SLR}} \\
\cmidrule(lr){2-4}\cmidrule(l){5-7}
\textbf{Grupo} & \textbf{DI} & \textbf{$t$ (gl)} & \textbf{$p$} & \textbf{DI} & \textbf{$t$ (gl)} & \textbf{$p$} \\
\midrule
Juvenil -- Control  & $0{,}242 \pm 0{,}050$ (9)  & 4,80 (8)  & \textbf{0,001} & $0{,}210 \pm 0{,}040$ (10) & 5,28 (9)  & \textbf{$<$\,0,001} \\
Juvenil -- Sacarosa & $-0{,}082 \pm 0{,}048$ (9) & $-1{,}70$ (8) & 0,127 & $-0{,}112 \pm 0{,}077$ (9) & $-1{,}46$ (8) & 0,183 \\
Juvenil -- Stevia   & $0{,}191 \pm 0{,}100$ (10) & 1,90 (9)  & 0,090 & $0{,}181 \pm 0{,}109$ (10) & 1,67 (9)  & 0,130 \\
Adulto -- Control   & $0{,}218 \pm 0{,}054$ (10) & 4,00 (9)  & \textbf{0,003} & $0{,}244 \pm 0{,}070$ (10) & 3,47 (9)  & \textbf{0,007} \\
Adulto -- Sacarosa  & $0{,}109 \pm 0{,}074$ (10) & 1,48 (9)  & 0,173 & $-0{,}082 \pm 0{,}110$ (10) & $-0{,}75$ (9) & 0,475 \\
Adulto -- Stevia    & $0{,}106 \pm 0{,}079$ (10) & 1,34 (9)  & 0,212 & $0{,}054 \pm 0{,}107$ (10) & 0,51 (9)  & 0,624 \\
\bottomrule
\end{tabularx}
\end{table}

\begin{table}[htbp]
\centering
\caption[Resumen del ANOVA de dos vías para el DI en la prueba SLR]{Resumen del ANOVA de dos vías (edad de exposición $\times$ tratamiento) para el DI en cada configuración de la prueba SLR. $\eta^2_p$: eta cuadrado parcial (tamaño del efecto).}
\label{tab:slr-anova}
\footnotesize
\begin{tabularx}{\linewidth}{@{}l C c c C c c@{}}
\toprule
& \multicolumn{3}{c}{\textbf{d-SLR}} & \multicolumn{3}{c}{\textbf{s-SLR}} \\
\cmidrule(lr){2-4}\cmidrule(l){5-7}
\textbf{Fuente de variación} & \textbf{$F$ (gl)} & \textbf{$p$} & \textbf{$\eta^2_p$} & \textbf{$F$ (gl)} & \textbf{$p$} & \textbf{$\eta^2_p$} \\
\midrule
Tratamiento & 4,58 (2, 52) & \textbf{0,015} & 0,150 & 6,67 (2, 53) & \textbf{0,003} & 0,201 \\
Edad        & 0,22 (1, 52) & 0,642 & 0,004 & 0,08 (1, 53) & 0,779 & 0,001 \\
Interacción & 2,04 (2, 52) & 0,140 & 0,073 & 0,53 (2, 53) & 0,591 & 0,020 \\
\bottomrule
\end{tabularx}
\end{table}

\begin{figure}[htbp]
  \centering
  \begin{subfigure}[b]{0.48\linewidth}
    \centering
    \incluirfigura[width=\linewidth]{Fig5_DI_SLR_d.pdf}
    \caption{d-SLR}
  \end{subfigure}
  \hfill
  \begin{subfigure}[b]{0.48\linewidth}
    \centering
    \incluirfigura[width=\linewidth]{Fig5_DI_SLR_s.pdf}
    \caption{s-SLR}
  \end{subfigure}
  \caption[Índice de discriminación (DI) en la prueba SLR]{Índice de discriminación (DI) en la prueba SLR. (A)~Configuración d-SLR. (B)~Configuración s-SLR. Las barras indican media $\pm$ SEM. Azul: control; rojo: sacarosa; verde: stevia (paleta común a todas las figuras de esta tesis). Se indican los resultados de las comparaciones de Tukey entre grupos (ns: no significativo; $^{*}\,p < 0{,}05$; $^{**}\,p < 0{,}01$). $n = 9$--$10$ animales por grupo.}
  \label{fig:slr-di}
\end{figure}

\begin{table}[htbp]
\centering
\caption[Tasa de exploración en la prueba SLR]{Tasa de exploración en la prueba SLR, equivalente al DI de la Tabla~\ref{tab:slr-di} ($\mathrm{DI} = 2\cdot\text{tasa} - 1$). Valores expresados como media $\pm$ SEM; entre paréntesis, número de animales.}
\label{tab:slr-tasa}
\small
\begin{tabular}{@{}lcc@{}}
\toprule
\textbf{Grupo} & \textbf{d-SLR} & \textbf{s-SLR} \\
\midrule
Juvenil -- Control  & $0{,}621 \pm 0{,}025$ (9)  & $0{,}605 \pm 0{,}020$ (10) \\
Juvenil -- Sacarosa & $0{,}459 \pm 0{,}024$ (9)  & $0{,}444 \pm 0{,}038$ (9) \\
Juvenil -- Stevia   & $0{,}595 \pm 0{,}050$ (10) & $0{,}590 \pm 0{,}054$ (10) \\
Adulto -- Control   & $0{,}609 \pm 0{,}027$ (10) & $0{,}622 \pm 0{,}035$ (10) \\
Adulto -- Sacarosa  & $0{,}554 \pm 0{,}037$ (10) & $0{,}459 \pm 0{,}055$ (10) \\
Adulto -- Stevia    & $0{,}553 \pm 0{,}039$ (10) & $0{,}527 \pm 0{,}053$ (10) \\
\bottomrule
\end{tabular}
\end{table}

\begin{figure}[htbp]
  \centering
  \begin{subfigure}[b]{0.48\linewidth}
    \centering
    \incluirfigura[width=\linewidth]{Fig5_tasa_SLR_d.pdf}
    \caption{d-SLR}
  \end{subfigure}
  \hfill
  \begin{subfigure}[b]{0.48\linewidth}
    \centering
    \incluirfigura[width=\linewidth]{Fig5_tasa_SLR_s.pdf}
    \caption{s-SLR}
  \end{subfigure}
  \caption[Tasa de exploración en la prueba SLR]{Tasa de exploración en la prueba SLR. (A)~Configuración d-SLR. (B)~Configuración s-SLR. Las barras indican media $\pm$ SEM. Azul: control; rojo: sacarosa; verde: stevia (paleta común a todas las figuras de esta tesis). Se indican los resultados de las comparaciones de Tukey entre grupos (ns: no significativo; $^{*}\,p < 0{,}05$; $^{**}\,p < 0{,}01$), estadísticamente idénticos a los del DI (Figura~\ref{fig:slr-di}) por tratarse de una transformación lineal. $n = 9$--$10$ animales por grupo.}
  \label{fig:slr-tasa}
\end{figure}

\FloatBarrier
\subsection{Síntesis}\label{sec:res-slr-sintesis}

En síntesis, los animales control discriminaron la posición novedosa en ambas configuraciones y en ambas edades de exposición. La exposición juvenil a sacarosa se asoció a un DI menor que el del control en ambas configuraciones, mientras que la exposición adulta a sacarosa lo hizo únicamente en la configuración s-SLR. Los grupos expuestos a stevia no difirieron del control en ninguna condición; en la configuración d-SLR, el grupo stevia juvenil presentó un DI mayor que el grupo sacarosa juvenil, aunque ninguno de los grupos stevia superó significativamente el nivel de azar. La Figura~\ref{fig:slr-tukey} resume las diferencias entre grupos junto con sus intervalos de confianza.

\begin{figure}[htbp]
  \centering
  \incluirfigura[width=\linewidth]{Fig5_estimacion_Tukey_SLR.png}
  \caption[Diferencias de medias del DI entre grupos (prueba de Tukey)]{Diferencias de medias del DI entre grupos (prueba de Tukey) con su intervalo de confianza del 95\,\%, para las configuraciones d-SLR (A) y s-SLR (B). Cada fila corresponde a una comparación; una diferencia positiva indica un DI mayor en el primer grupo de la comparación. En rojo se indican las comparaciones significativas ($p < 0{,}05$). La línea punteada indica ausencia de diferencia.}
  \label{fig:slr-tukey}
\end{figure}

% ---------------------------------------------------------------------
\FloatBarrier
\section{Neurogénesis hipocampal (PCNA y DCX)}\label{sec:res-neurogenesis}
\pendiente{redactar cuando esté el análisis histológico.}

\FloatBarrier
\section{Análisis multivariado y aprendizaje automático}\label{sec:res-ml}
\pendiente{PCA, clustering jerárquico (Ward) y Random Forest (objetivos M1 y M2).}
